\begin{table}[htbp]
\centering
\caption{Border-vs-Interior Estimates: Revenue per Store, Pooled Colorado + Washington + Oregon}
\label{tab:pooled_perstore_or}
\begin{tabular}{lcccc}
\toprule
Specification & Coefficient & SE (HC1) & SE (county-cl.) & $p$-value (county-cl.) \\
\midrule
No controls & $+0.526$ & $0.099$ & $0.226$ & $0.020$ \\
+ log(pop) + year FE + state FE & $+0.429$ & $0.074$ & $0.157$ & $0.006$ \\
\quad excluding Sedgwick, CO and Malheur, OR & $+0.323$ & $0.067$ & $0.142$ & $0.023$ \\
\midrule
\multicolumn{5}{l}{\textit{Single-state estimates, + log(pop) [+ year FE for CO/WA] (all HC1 $p<0.01$)}} \\
Colorado only & $+0.552$ & $0.088$ & -- & -- \\
Washington only & $+0.441$ & $0.101$ & -- & -- \\
Oregon only (2025 cross-section) & $+1.168$ & $0.433$ & -- & -- \\
\bottomrule
\end{tabular}
\begin{minipage}{0.95\textwidth}
\vspace{0.3em}
\footnotesize
\textit{Notes:} Adds Oregon to the Colorado+Washington revenue-per-store panel in
Table~\ref{tab:revenue_perstore_border70}. Oregon's county sales come from OLCC's own Socrata series
(2017--2025, all years complete); Oregon's store counts, unlike Colorado's and Washington's, cannot be
built as a historical cumulative-open panel -- the geocoded dispensary listing used elsewhere in this
project has real coverage gaps for Oregon specifically (several records geocode outside Oregon entirely,
and Baker County, one of the three border counties, has zero matches despite recording revenue in every
year 2017--2025) -- so Oregon instead contributes a single 2025 cross-section, pairing 2025 revenue with
currently-active retailer counts from OLCC's own license registry (which tags county natively). Pooling
this cross-section into the panel is valid despite the year mismatch: year fixed effects absorb any
Oregon-specific 2025 level shock the same way they absorb Colorado's and Washington's own year effects,
and state fixed effects absorb any state-level difference in revenue per store, so Oregon's within-state
border/interior contrast is identified the same way Colorado's and Washington's are. Adding Oregon
\emph{strengthens} the pooled estimate's significance under county-clustered SEs ($p=0.015\to p=0.006$),
and all three states are independently positive and significant. Malheur County, OR (Ontario, directly
across the Snake River from Fruitland/Payette, Idaho) plays the same outlier role for Oregon that Sedgwick
County plays for Colorado; excluding both, the pooled estimate remains significant.
\end{minipage}
\end{table}
